%! AP Statistics — Unit 1, Lesson 1.2 — Guided Notes ANSWER KEY
\documentclass[10pt]{article}
\usepackage{apstats-article}
\usepackage{apstats-key}

%=============================================================================
\begin{document}

\pageheader{Unit 1, Lesson 1.2}{Guided Notes --- Answer Key}
\begin{tabularx}{\linewidth}{X X X}
Name:\;\ans{Key} & Date:\; & Period:\;
\end{tabularx}

\vspace{0.12in}

\begin{objectivebox}
\small By the end of this lesson, I will be able to\dots\\[4pt]
\ansline{classify a variable as categorical or quantitative and justify using the average test.}\\
\ansline{distinguish between discrete and continuous quantitative variables with examples.}\\
\ansline{explain why variable type must be identified before choosing a display or summary.}
\end{objectivebox}

\vspace{0.1in}

\begin{vocabbox}
\small
\begin{multicols}{2}
\noindent\textbf{\textcolor{navy}{Variable:}}\\
\ans{A characteristic that can take different values for different individuals}\\[6pt]

\noindent\textbf{\textcolor{navy}{Categorical variable:}}\\
\ans{Places an individual into one of several groups or categories; no meaningful arithmetic on the values}\\[6pt]

\noindent\textbf{\textcolor{navy}{Quantitative variable:}}\\
\ans{Takes numerical values for which arithmetic (average, difference) makes sense}\\[6pt]

\columnbreak

\noindent\textbf{\textcolor{navy}{Discrete quantitative variable:}}\\
\ans{Countable values with gaps between them (e.g., number of siblings: 0, 1, 2, \ldots)}\\[6pt]

\noindent\textbf{\textcolor{navy}{Continuous quantitative variable:}}\\
\ans{Can take any value in an interval; measured, not counted (e.g., height, time)}\\[6pt]

\noindent\textbf{\textcolor{navy}{Distribution of a variable:}}\\
\ans{Describes what values the variable takes and how often each value occurs}\\[6pt]
\end{multicols}
\end{vocabbox}

\vspace{0.1in}

\begin{hookbox}
\small\textbf{Student Data Table}

\vspace{8pt}
\renewcommand{\arraystretch}{1.5}
\begin{tabularx}{\linewidth}{@{} C{0.18\linewidth} C{0.14\linewidth}
                                    C{0.14\linewidth} C{0.2\linewidth}
                                    C{0.2\linewidth} @{}}
\rowcolor{sky}
\textbf{Name} & \textbf{Grade} & \textbf{GPA} &
\textbf{Fav.\ Color} & \textbf{Sleep (hrs)} \\
\midrule
Jordan  & 11 & 3.7 & Blue   & 7.5 \\
Priya   & 10 & 3.2 & Green  & 6.0 \\
Marcus  & 12 & 2.9 & Red    & 8.0 \\
Sofia   & 11 & 3.5 & Purple & 7.0 \\
\end{tabularx}

\vspace{10pt}

\renewcommand{\arraystretch}{1.8}
\begin{tabularx}{\linewidth}{@{} |>{\centering\arraybackslash}X
                                    |>{\centering\arraybackslash}X| @{}}
\hline
\textbf{Group A (Categorical)} & \textbf{Group B (Quantitative)} \\
\hline
\ans{Name, Favorite Color} & \ans{Grade level, GPA, Sleep (hrs)} \\
\hline
\end{tabularx}

\vspace{8pt}
\textbf{Sorting rule:}\quad
\ans{Group A values are words/labels; Group B values are numbers you can average meaningfully.}

\begin{teachernote}
Grade level is quantitative (discrete) --- you can average it even if it's rarely useful to do so.
Favorite Color is clearly categorical. Name is a label, not a variable of interest for analysis.
Accept students grouping Name with categorical.
\end{teachernote}
\end{hookbox}

\vspace{0.1in}

\begin{notesbox}{Classifying Variables}
\small
\begin{multicols}{2}

\textbf{\textcolor{navy}{The Two Main Types}}\\[4pt]
\textbf{Categorical variable:}\\
Places an individual into one of several \ans{groups or categories}.
No meaningful arithmetic.\\[6pt]
\textit{Examples:}
\begin{itemize}[itemsep=2pt]
  \item Favorite color: \ans{categorical}
  \item Mode of transport: \ans{categorical}
\end{itemize}

\vspace{8pt}
\textbf{Quantitative variable:}\\
Takes numerical values for which \ans{arithmetic} makes sense.\\[6pt]
\textit{Examples:}
\begin{itemize}[itemsep=2pt]
  \item GPA: \ans{quantitative}
  \item Hours of sleep: \ans{quantitative}
\end{itemize}

\columnbreak

\textbf{\textcolor{navy}{The Average Test}}\\[4pt]
\textit{``Does averaging these values give a \ans{meaningful} result?''}\\[6pt]
\begin{tabularx}{\linewidth}{@{} l X @{}}
\textbf{Yes} $\to$ & \ans{Quantitative} \\[4pt]
\textbf{No}  $\to$ & \ans{Categorical} \\
\end{tabularx}

\vspace{10pt}
\textbf{\textcolor{navy}{Trap Variables}}\\[4pt]
\begin{itemize}[itemsep=4pt]
  \item ZIP code: \ans{categorical}
  \item Jersey number: \ans{categorical}
  \item Student ID: \ans{categorical}
  \item Likert rating (1--5): \ans{debatable}
\end{itemize}
\end{multicols}
\end{notesbox}

\vspace{0.1in}

\begin{notesbox}{Discrete vs.\ Continuous Quantitative Variables}
\small
\begin{multicols}{2}

\textbf{\textcolor{navy}{Discrete}}\\
Values are \ans{countable} --- there are \ans{gaps} between possible values.\\[4pt]
Ask: \textit{``Can I \ans{count} them?''}\\[6pt]
\textit{Example:} Number of pets\\
Possible values: \ans{0, 1, 2, 3, \ldots}\\[4pt]
\begin{itemize}[itemsep=2pt]
  \item Number of siblings: \ans{discrete}
  \item Number of songs: \ans{discrete}
\end{itemize}

\columnbreak

\textbf{\textcolor{navy}{Continuous}}\\
Can take \ans{any} value in an interval --- \ans{measured}, not counted.\\[4pt]
Ask: \textit{``Can it fall \ans{between} two values?''}\\[6pt]
\textit{Example:} Weight of a pet\\
Possible values: \ans{any positive real number}\\[4pt]
\begin{itemize}[itemsep=2pt]
  \item Height in cm: \ans{continuous}
  \item Commute time: \ans{continuous}
\end{itemize}
\end{multicols}

\vspace{6pt}
\renewcommand{\arraystretch}{1.5}
\begin{tabularx}{\linewidth}{@{} >{\centering\arraybackslash\columncolor{navy}}p{0.28\linewidth}
                                    >{\centering\arraybackslash}X
                                    >{\centering\arraybackslash}X @{}}
\textcolor{white}{\textbf{Question}} &
\textcolor{navy}{\textbf{Discrete}} &
\textcolor{navy}{\textbf{Continuous}} \\
Can I list all possible values? & \ans{Yes} & \ans{No} \\
Is it measured on a scale?      & \ans{No}  & \ans{Yes} \\
Are there gaps between values?  & \ans{Yes} & \ans{No} \\
\end{tabularx}
\end{notesbox}

\vspace{0.1in}

\begin{practicebox}
\small\textbf{Classify each variable.}

\vspace{6pt}
\renewcommand{\arraystretch}{2.0}
\begin{tabularx}{\linewidth}{@{} c >{\raggedright\arraybackslash}p{0.38\linewidth}
                                    C{0.1\linewidth} X @{}}
\rowcolor{sky}
\# & \textbf{Variable} & \textbf{Type} & \textbf{Justification} \\
\midrule
1 & Eye color              & \ans{C}  & \ans{Labels a category; averaging eye colors is meaningless} \\
2 & Number of siblings     & \ans{QD} & \ans{Countable whole numbers; gaps between values} \\
3 & Body temperature (°F)  & \ans{QC} & \ans{Measured on a continuous scale; can be any decimal} \\
4 & Area code              & \ans{C}  & \ans{Numeric label for a geographic region; averaging gives no meaning} \\
5 & Number of words in essay & \ans{QD} & \ans{Count of whole words; gaps between values} \\
6 & Time to run a mile     & \ans{QC} & \ans{Measured time; can take any positive value in an interval} \\
\end{tabularx}
\end{practicebox}

\vspace{0.1in}

\begin{spiralbox}
\small
\begin{multicols}{2}

\textbf{This lesson connects forward to\dots}
\begin{itemize}[itemsep=2pt]
  \item \textbf{Lessons 1.3--1.4:} Displaying categorical data (bar charts, two-way tables)
  \item \textbf{Lessons 1.5--1.9:} Displaying and summarizing quantitative data
  \item \textbf{Unit 4:} Discrete probability distributions vs.\ density curves
\end{itemize}

\columnbreak

\textbf{Big Idea:}\\[4pt]
\ans{Variable type determines which tools are valid. Using a mean for a categorical
variable (like averaging ZIP codes) produces a number that means nothing. Picking the
wrong tool leads to wrong conclusions --- so identifying variable type is always Step 0
in data analysis.}

\end{multicols}
\end{spiralbox}

\end{document}
